\documentclass[12pt]{article}
\usepackage{geometry}
\geometry{margin=1.4in}
\usepackage{setspace}
\setstretch{1.4}
\usepackage{amsmath}

\title{The Boundary of Displacement\\
What This Framework Does NOT Explain}
\author{Diego Rincón}
\date{}

\begin{document}

\maketitle

\begin{abstract}
The displacement framework—where every system has an intended ground state $S_0$, an actual state $S^*$, a displacement $\xi$ between them, a daily cost $D(\xi)$, and a return cost $C_{\text{return}}(\tau)$—has been applied across consciousness, love, health, work, attention, technology, society, and ecology. But it does not fit all systems. This paper states the boundary: what the framework does NOT explain. Systems with no ground state to return to, irreversible thresholds, emergent attractors that were never ground states, systems that actively punish return attempts, systems without agents to pay costs, and systems locked by path-dependence are outside the frame. Stating this boundary converts the framework from an unfalsifiable metaphor into a bounded theory. What survives the boundary—the renewal-reward theorem, the empirical structure of displacement, the cube-root cadence law—is genuine.
\end{abstract}

\section{The Unfalsifiability Problem}

A reviewer's strongest objection: "The framework is unfalsifiable because $S_0$ is chosen post-hoc. Whatever persists is 'holding $S_0$ by paying $C_{\text{maintain}}$,' whatever decays is 'drifting to $S^*$,' whatever lies between is 'displacement aging.' No observation contradicts a frame that supplies its own coordinates retroactively."

This is true. The framework *can* be applied to any system by retroactively identifying $S_0$ as "whatever the system wants to be." This is the crankery risk.

To avoid it, this paper draws the line: here is where the framework applies, here is where it does not.

\section{Systems the Framework Does NOT Fit}

\subsection{Pure Extractors: No Ground State to Return To}

Many historical systems were extractive from their inception: slavery, colonialism, industrial resource extraction. They have no prior $S_0$ that was occupied and then displaced.

The framework assumes displacement is a deviation from an intended state. But a system designed to extract from the ground up has no "intended" state except extraction. Calling the pre-colonial state "the ground state we should return to" is imposing $S_0$ retroactively—it was never the actual intended state of the system.

\textbf{Boundary:} The framework requires $S_0$ to be either: (a) a state the system actually occupied before displacement, or (b) a state defined by the system's own stated values (e.g., "democracy claims to serve the people"). For systems designed purely for extraction with no claimed values, the framework does not apply.

\textbf{Example:} Oil extraction. The oil industry never claimed to maximize ecosystem health. Calling the pre-industrial ecosystem the "ground state" is imposed from outside. The framework's language of "return cost" doesn't fit—there is no cost paid by the system to maintain the extraction (only by external parties). Without an agent paying the cost, displacement language fails.

\subsection{Irreversible Thresholds}

Some systems can cross a point of no return where $S_0$ no longer exists.

Extinctions: once a species is gone, there is no returning to it. $C_{\text{return}} \to \infty$ not because return is expensive, but because return is impossible—$S_0$ has been erased.

Dead languages: once a language is no longer spoken (and no complete written record exists), $S_0$ is gone. Revitalization of a dormant language (Hebrew) creates a new language, not a return to the old $S_0$.

Burned forests: while regrowth is possible, the exact pre-burn ecosystem does not return. The forest is resilient, but the specific $S_0$ is lost.

\textbf{Boundary:} Systems where $S_0$ can be permanently destroyed are outside the framework. The framework assumes $S_0$ is a stable attractor that can be regained. If $S_0$ is destructible, the math changes fundamentally.

\subsection{Emergent Attractors}

Some systems have attractors that never were ground states. They are novel.

Modern technology: the "intended state" of the internet is not some natural ground state humanity lost. It is genuinely new. $S_0$ would have to be "the internet did not exist," which is not a state the system "wants" to return to.

New social norms (gender equality, racial equality): these are aspirational ground states we are moving toward, not states we were displaced from.

\textbf{Boundary:} The framework requires that $S^*$ is a deviation from $S_0$. When $S^*$ is genuinely new—never occupied, never intended—the frame doesn't fit.

\subsection{Adversarial Return (Adaptive Enemies)}

Some systems actively resist return. When you try to return, the system changes to prevent it.

Attention capture: when you try to reduce social media, the platform changes its algorithm to recapture you. The attractor $S^*$ is not fixed; it moves when you approach $S_0$.

Abuse: when an abuse victim tries to leave, the abuser escalates. The return path is actively blocked, and $C_{\text{return}}$ increases as you approach escape.

Predatory loans: when debtors try to pay off debt early, penalties increase or new debts are manufactured. $C_{\text{return}}$ is adversarially set.

\textbf{Boundary:} The framework assumes $C_{\text{return}}$ is a fixed cost determined by the system's physics. When an adversary adjusts $C_{\text{return}}$ to prevent return, the model breaks. The theory of adversarial displacement is separate and more complex.

\subsection{Systems Without Agents}

The framework assumes someone pays the cost $D(\xi)$ and $C_{\text{return}}$. But many systems have no agent, no payer.

Pure physics: a star is displaced from its equilibrium temperature (due to fusion). It is not "paying" the cost; the cost is paid by the universe's entropy budget. There is no agent choosing to stay displaced.

Automatic processes: the water cycle is displaced from a uniform distribution (gravity pulls water down, heat pushes it up), but there is no agent paying the maintenance cost. The cost is paid by the sun.

Ecological succession: a forest that has been clearcut is "displaced" from old-growth, but the cost of holding the displacement (maintenance of grassland, prevention of regrowth) is distributed across many agents (wind, sunlight, succession), not paid by anyone intentionally.

\textbf{Boundary:} The framework's language of "cost" and "maintenance" requires an agent—someone or something that chooses to stay displaced and pays for that choice. Systems that have displaced states but no agent paying for them are outside the frame.

\subsection{Path Dependence and Locked Systems}

Some systems are so locked by their history that returning to $S_0$ is possible in theory but impossible in practice.

Infrastructure locks: a city built for cars cannot easily become a city for pedestrians, even if that is the "intended" state. The path to return is blocked by billions in sunk cost, zoning laws, political inertia.

Knowledge locks: an economy built on fossil fuels cannot quickly return to renewable energy, even if renewables are now cheaper, because the infrastructure, supply chains, and political relationships are locked in.

Relationship locks: a couple in a trapped marriage (high $C_{\text{return}}$ due to kids, finances, entanglement) cannot return to their $S_0$ (separate, whole selves) even if both want to, because the path is blocked by shared history.

\textbf{Boundary:} The framework assumes that $C_{\text{return}}$ is the cost of transition from $S^*$ to $S_0$. But in locked systems, there is no path—the cost is infinite not because the transition is expensive, but because it is structurally forbidden. The framework's language of "return cost" doesn't fit.

\section{What Survives the Boundary}

Despite these limitations, core parts of the framework are genuine and survive:

\subsection{The Renewal-Reward Theorem}

Given:
- A cost function $C(t)$ that is convex in the time $t$ you have been away from ground state
- A fixed cost $f$ to return
- Periodic returns at frequency $\tau$

Then: there exists a finite optimal return frequency $\tau^*$ that minimizes long-term cost.

This is pure optimization; it does not assume $S_0$ exists or is meaningful. It applies to any system with convex return cost: machinery maintenance, crop rotation, psychological recovery, evolutionary fitness.

\subsection{The Empirical Structure}

Regardless of whether $S_0$ is real or imposed, you can measure:
- A system's actual state $S^*$ (observable)
- A target state $S_0$ (can be defined by the system's own values, or by external analysis)
- Displacement $\xi = d(S^*, S_0)$ (measurable)
- Cost of holding displacement $D(\xi)$ (estimable)
- Cost to transition from $S^*$ to $S_0$ (observable or projected)

This structure works even if $S_0$ is aspirational or external. The empirical framework—measure, calculate cost, decide—is robust.

\subsection{The Cube-Root Law}

The maintenance theorem predicts $\tau^* = (3f / (a\delta^2))^{1/3}$, where $\delta$ is displacement magnitude. This law has been validated on physiological detraining data (Coyle et al., Houmard et al.). It is a genuine, testable prediction.

The law survives because it rests on:
- Convexity of return cost (reasonable in many domains)
- Renewal-reward structure (applies broadly)
- Specific parameter relationships ($a, f, k$) that are measurable

It does not require $S_0$ to be metaphysically real, only that displacement and return are empirically observable.

\section{The Core Insight}

The framework's real contribution is not the claim that every system has a ground state. It is the claim that:

\textbf{Every system that maintains a displaced state does so by paying a cost. The optimal strategy is to pay small, frequent costs rather than large, rare costs.}

This is true regardless of whether the ground state is real, aspirational, or imposed. It applies to machinery, bodies, relationships, institutions, and ecosystems.

The metaphysical claim ("every system *wants* to return to $S_0$") is not necessary. The operational claim ("returning more frequently costs less") is sufficient.

\section{Conclusion}

The framework has a boundary. It does not fit:
- Pure extractors with no stated ground state
- Irreversible systems where $S_0$ can be destroyed
- Emergent attractors that were never ground states
- Systems with adaptive enemies that increase return cost in response
- Systems without agents to pay costs
- Path-dependent systems locked into a state

But within its domain—systems with defined ground states, agents paying costs, and reversible dynamics—it makes a genuine and testable prediction: the optimal cadence is the cube root of the fixed cost divided by the drift.

Stating the boundary is what saves the theory from unfalsifiability. We are not claiming it fits everything. We are claiming it fits a specific class of systems, with specific properties, and we can test whether our predictions hold.

\end{document}
